\documentclass[10pt,a4paper]{amsart}
\usepackage{amsmath,amssymb,amsthm,graphicx}
% Retained minimal preamble for the Stacks Project chapter more-morphisms.tex (More on Morphisms).
% Verbatim excerpt of the chapter preamble: Xy-pic setup (native preamble lines 14-18) and the native
% theorem declarations with the mathematical macros (native preamble lines 154-217). Copied mechanically;
% nothing added, removed or redefined. The wrapper class and the amsmath/amssymb/amsthm/graphicx base
% packages are supplied by the assembler, so the class/externaldocument part of the native preamble is omitted.
% For commutative diagrams we use Xy-pic
\usepackage[all]{xy}

% We use 2cell for 2-commutative diagrams.
\xyoption{2cell}
% Theorem environments.
%
\theoremstyle{plain}
\newtheorem{theorem}[subsection]{Theorem}
\newtheorem{proposition}[subsection]{Proposition}
\newtheorem{lemma}[subsection]{Lemma}

\theoremstyle{definition}
\newtheorem{definition}[subsection]{Definition}
\newtheorem{example}[subsection]{Example}
\newtheorem{exercise}[subsection]{Exercise}
\newtheorem{situation}[subsection]{Situation}

\theoremstyle{remark}
\newtheorem{remark}[subsection]{Remark}
\newtheorem{remarks}[subsection]{Remarks}

\numberwithin{equation}{subsection}

% Macros
%
\def\lim{\mathop{\mathrm{lim}}\nolimits}
\def\colim{\mathop{\mathrm{colim}}\nolimits}
\def\Spec{\mathop{\mathrm{Spec}}}
\def\Hom{\mathop{\mathrm{Hom}}\nolimits}
\def\Ext{\mathop{\mathrm{Ext}}\nolimits}
\def\SheafHom{\mathop{\mathcal{H}\!\mathit{om}}\nolimits}
\def\SheafExt{\mathop{\mathcal{E}\!\mathit{xt}}\nolimits}
\def\Sch{\mathit{Sch}}
\def\Mor{\mathop{\mathrm{Mor}}\nolimits}
\def\Ob{\mathop{\mathrm{Ob}}\nolimits}
\def\Sh{\mathop{\mathit{Sh}}\nolimits}
\def\NL{\mathop{N\!L}\nolimits}
\def\CH{\mathop{\mathrm{CH}}\nolimits}
\def\proetale{{pro\text{-}\acute{e}tale}}
\def\etale{{\acute{e}tale}}
\def\QCoh{\mathit{QCoh}}
\def\Ker{\mathop{\mathrm{Ker}}}
\def\Im{\mathop{\mathrm{Im}}}
\def\Coker{\mathop{\mathrm{Coker}}}
\def\Coim{\mathop{\mathrm{Coim}}}

% Boxtimes
%
\DeclareMathSymbol{\boxtimes}{\mathbin}{AMSa}{"02}

%
% Macros for moduli stacks/spaces
%
\def\QCohstack{\mathcal{QC}\!\mathit{oh}}
\def\Cohstack{\mathcal{C}\!\mathit{oh}}
\def\Spacesstack{\mathcal{S}\!\mathit{paces}}
\def\Quotfunctor{\mathrm{Quot}}
\def\Hilbfunctor{\mathrm{Hilb}}
\def\Curvesstack{\mathcal{C}\!\mathit{urves}}
\def\Polarizedstack{\mathcal{P}\!\mathit{olarized}}
\def\Complexesstack{\mathcal{C}\!\mathit{omplexes}}
% \Pic is the operator that assigns to X its picard group, usage \Pic(X)
% \Picardstack_{X/B} denotes the Picard stack of X over B
% \Picardfunctor_{X/B} denotes the Picard functor of X over B
\def\Pic{\mathop{\mathrm{Pic}}\nolimits}
\def\Picardstack{\mathcal{P}\!\mathit{ic}}
\def\Picardfunctor{\mathrm{Pic}}
\def\Deformationcategory{\mathcal{D}\!\mathit{ef}}
\newenvironment{reference}{}{}
\newenvironment{slogan}{}{}
\newenvironment{history}{}{}
\usepackage{cleveref}
\crefname{originallemma}{original lemma}{original lemmas}
\Crefname{originallemma}{Original lemma}{Original lemmas}
\makeatletter
\newlabel{algebra-lemma-surjective-endo-noetherian-ring-is-iso}{{algebra-lemma-surjective-endo-noetherian-ring-is-iso}{source}{}{sourcepointer-1}{}}
\newlabel{algebra-lemma-surjective-endo-noetherian-ring-is-iso@cref}{{[originallemma][0][]algebra-lemma-surjective-endo-noetherian-ring-is-iso}{[1][0][]source}{}{}{}}
\makeatother
\crefname{originallemma}{original lemma}{original lemmas}
\Crefname{originallemma}{Original lemma}{Original lemmas}
\makeatletter
\newlabel{lemma-control-graded}{{lemma-control-graded}{source}{}{sourcepointer-2}{}}
\newlabel{lemma-control-graded@cref}{{[originallemma][0][]lemma-control-graded}{[1][0][]source}{}{}{}}
\makeatother
\crefname{originallemma}{original lemma}{original lemmas}
\Crefname{originallemma}{Original lemma}{Original lemmas}
\makeatletter
\newlabel{lemma-lifting-along-artinian-at-point}{{lemma-lifting-along-artinian-at-point}{source}{}{sourcepointer-3}{}}
\newlabel{lemma-lifting-along-artinian-at-point@cref}{{[originallemma][0][]lemma-lifting-along-artinian-at-point}{[1][0][]source}{}{}{}}
\makeatother
\crefname{originallemma}{original lemma}{original lemmas}
\Crefname{originallemma}{Original lemma}{Original lemmas}
\makeatletter
\newlabel{lemma-relative-map-approximation}{{lemma-relative-map-approximation}{source}{}{sourcepointer-4}{}}
\newlabel{lemma-relative-map-approximation@cref}{{[originallemma][0][]lemma-relative-map-approximation}{[1][0][]source}{}{}{}}
\makeatother
\crefname{originallemma}{original lemma}{original lemmas}
\Crefname{originallemma}{Original lemma}{Original lemmas}
\makeatletter
\newlabel{more-algebra-lemma-flat-unramified}{{more-algebra-lemma-flat-unramified}{source}{}{sourcepointer-5}{}}
\newlabel{more-algebra-lemma-flat-unramified@cref}{{[originallemma][0][]more-algebra-lemma-flat-unramified}{[1][0][]source}{}{}{}}
\makeatother
\crefname{originallemma}{original lemma}{original lemmas}
\Crefname{originallemma}{Original lemma}{Original lemmas}
\makeatletter
\newlabel{morphisms-lemma-etale-smooth-unramified}{{morphisms-lemma-etale-smooth-unramified}{source}{}{sourcepointer-6}{}}
\newlabel{morphisms-lemma-etale-smooth-unramified@cref}{{[originallemma][0][]morphisms-lemma-etale-smooth-unramified}{[1][0][]source}{}{}{}}
\makeatother
\title{Approximating a relative map that is an isomorphism on completions by an etale morphism}
\author{The Stacks Project Authors}
\date{stacks-project:more-morphisms.tex:local-snapshot-20261011; GNU FDL 1.2 or later}
\begin{document}
\maketitle
\noindent\textit{Editorial note.} Statement span is lines 11720-11725 of sources/stacks-more-morphisms/pinned-original-source.tex; proof span is lines 11727-11798, the complete author proof (70 body lines).\par\medskip
\noindent\textit{Editorial note.} Navigation stubs: 6 cited-but-not-included labels are supplied as non-load-bearing pointers with their pinned tags where the snapshot tag index has them.\par\medskip
\noindent\textit{Editorial note.} extra\_packages is empty on purpose: the retained preamble already loads xy with options and a second declaration clashes in the assembler wrapper.\par\medskip
\noindent\textit{Editorial note.} Text is retained verbatim, including the author's own typos and placeholders.\par\medskip
\noindent\textit{Allowed standard prerequisites (editorial).}\begin{itemize}
\item Explicit prerequisite boundary. The relative map approximation lemma, completion of local rings, elementary etale neighbourhoods, and the lifting criterion for smoothness; the reduction is retained.
\end{itemize}\medskip
\ifdefined\hypertarget\hypertarget{sourcepointer-1}{}\fi
\noindent\textit{Source reference.} Original lemma algebra-lemma-surjective-endo-noetherian-ring-is-iso: sources/stacks-more-morphisms/pinned-original-source.tex, source lines 11798--11798. Cited by the retained statement or proof but not part of the collected spans; supplied only as a navigation pointer. The argument does not depend on its content.\par\medskip
\ifdefined\hypertarget\hypertarget{sourcepointer-2}{}\fi
\noindent\textit{Source reference.} Original lemma lemma-control-graded: sources/stacks-more-morphisms/pinned-original-source.tex, source lines 11571--11590. Cited by the retained statement or proof but not part of the collected spans; supplied only as a navigation pointer. The argument does not depend on its content.\par\medskip
\ifdefined\hypertarget\hypertarget{sourcepointer-3}{}\fi
\noindent\textit{Source reference.} Original lemma lemma-lifting-along-artinian-at-point: sources/stacks-more-morphisms/pinned-original-source.tex, source lines 3575--3600. Cited by the retained statement or proof but not part of the collected spans; supplied only as a navigation pointer. The argument does not depend on its content.\par\medskip
\ifdefined\hypertarget\hypertarget{sourcepointer-4}{}\fi
\noindent\textit{Source reference.} Original lemma lemma-relative-map-approximation: sources/stacks-more-morphisms/pinned-original-source.tex, source lines 11289--11363. Cited by the retained statement or proof but not part of the collected spans; supplied only as a navigation pointer. The argument does not depend on its content.\par\medskip
\ifdefined\hypertarget\hypertarget{sourcepointer-5}{}\fi
\noindent\textit{Source reference.} Original lemma more-algebra-lemma-flat-unramified: sources/stacks-more-morphisms/pinned-original-source.tex, source lines 11798--11798. Cited by the retained statement or proof but not part of the collected spans; supplied only as a navigation pointer. The argument does not depend on its content.\par\medskip
\ifdefined\hypertarget\hypertarget{sourcepointer-6}{}\fi
\noindent\textit{Source reference.} Original lemma morphisms-lemma-etale-smooth-unramified: sources/stacks-more-morphisms/pinned-original-source.tex, source lines 11798--11798. Cited by the retained statement or proof but not part of the collected spans; supplied only as a navigation pointer. The argument does not depend on its content.\par\medskip
\begin{lemma}
\label{lemma-relative-isomorphism-approximation}
With notation an assumptions as in Lemma \ref{lemma-relative-map-approximation}
assume that $\varphi$ induces an isomorphism on completions.
Then we can choose our diagram such that $f$ is \'etale.
\end{lemma}

\begin{proof}
We may assume $N \geq 2$ and we may replace $(T, t)$ with $(T', t')$ as in
Lemma \ref{lemma-control-graded}. Since $(V, v) \to (S, s)$ is an elementary
\'etale neighbourhood, so is $(X \times_S V, (x, v)) \to (X, x)$.
Thus $\mathcal{O}_{X, x} \to \mathcal{O}_{X \times_S V, (x, v)}$
induces an isomorphism on completions by
More on Algebra, Lemma \ref{more-algebra-lemma-flat-unramified}.
We claim $\mathcal{O}_{X, x} \to \mathcal{O}_{W, w}$ induces
an isomorphism on completions. Having proved this,
Lemma \ref{lemma-lifting-along-artinian-at-point}
will show that $f$ is smooth at $w$ and of course
$f$ is unramified at $u$ as well, so
Morphisms, Lemma \ref{morphisms-lemma-etale-smooth-unramified}
tells us $f$ is \'etale at $w$.

\medskip\noindent
First we use the commutativity in part (5) of
Lemma \ref{lemma-relative-map-approximation}
to see that for $i \leq N$ there is a commutative diagram
$$
\xymatrix{
\text{Gr}^i_{\mathfrak m_x}(\mathcal{O}_{X, x})
\ar[r]_-\varphi &
\text{Gr}^i_{\mathfrak m_{y_\sigma}}(\mathcal{O}_{Y_\sigma, y_\sigma}^\wedge)
\ar@{=}[r] &
\text{Gr}^i_{\mathfrak m_{(y, v)}}(\mathcal{O}_{Y \times_{T, \tau} V, (y, v)})
\ar[d]_{\cong} \\
\text{Gr}^i_{\mathfrak m_x}(\mathcal{O}_{X, x})
\ar[r]^-{\cong} \ar@{=}[u] &
\text{Gr}^i_{\mathfrak m_{(x, v)}}(\mathcal{O}_{X \times_S V, (x, v)})
\ar[r]^{f^\sharp_w} &
\text{Gr}^i_{\mathfrak m_w}(\mathcal{O}_{W, w})
}
$$
This implies that $f^\sharp_w$ defines an isomorphism
$\kappa(x) \to \kappa(w)$ on residue fields and an isomorphism
$\mathfrak m_x/\mathfrak m_x^2 \to \mathfrak m_w/\mathfrak m_w^2$
on cotangent spaces. Hence $f^\sharp_w$ defines a surjection
$\mathcal{O}_{X, x}^\wedge \to \mathcal{O}_{W, w}^\wedge$
on complete local rings.

\medskip\noindent
By Lemma \ref{lemma-control-graded} there is an isomorphism of
$\text{Gr}_{\mathfrak m_s}(\mathcal{O}_{(Y \times_{T, \tau} V, (y, v)})$
with
$\text{Gr}_{\mathfrak m_s}(\mathcal{O}_{Y_\sigma, y_\sigma})$.
This follows by taking stalks of the isomorphism of conormal
sheaves at the point $y$. Since our local rings are Noetherian
taking associated graded with respect to $\mathfrak m_s$
commutes with completion because completion with respect to an ideal
is an exact functor on finite modules over Noetherian rings.
This produces the right vertical isomorphism in the diagram of graded rings
$$
\xymatrix{
\text{Gr}_{\mathfrak m_s}(\mathcal{O}_{W, w}^\wedge) &
\text{Gr}_{\mathfrak m_s}
(\mathcal{O}_{(Y \times_{T, \tau} V, (y, v)}^\wedge) \ar[l] \\
\text{Gr}_{\mathfrak m_s}(\mathcal{O}_{X, x}^\wedge)
\ar[r]^\varphi \ar[u] &
\text{Gr}_{\mathfrak m_s}(\mathcal{O}_{Y_\sigma, y_\sigma}^\wedge)
\ar[u]_{\cong}
}
$$
We do not claim the diagram commutes. By the result of the previous
paragraph the left arrow is surjective. The other three arrows
are isomorphisms. It follows that the left arrow is a surjective map
between isomorphic Noetherian rings. Hence it is an isomorphism
by Algebra, Lemma \ref{algebra-lemma-surjective-endo-noetherian-ring-is-iso}
(you can argue this directly using Hilbert functions as well).
In particular $\mathcal{O}_{X, x}^\wedge \to \mathcal{O}_{W, w}^\wedge$
must be injective as well as surjective which finishes the proof.
\end{proof}

\end{document}
